\documentclass[11pt,a4paper]{article}
\usepackage{fontspec}
\setmainfont{TeX Gyre Pagella}
\newfontfamily\headfont{Inter}[UprightFont=*-ExtraBold,BoldFont=*-ExtraBold,Extension=.otf,Path=/usr/share/fonts/opentype/inter/]
\usepackage[a4paper,margin=2.4cm]{geometry}
\usepackage{graphicx,booktabs,float,caption,array,xcolor,longtable,threeparttable}
\usepackage{xurl}
\usepackage[hidelinks]{hyperref}
\usepackage{titlesec}
\usepackage[round]{natbib}
\usepackage{setspace}
\setstretch{1.08}
\renewcommand{\topfraction}{0.9}\renewcommand{\bottomfraction}{0.8}\renewcommand{\textfraction}{0.1}\renewcommand{\floatpagefraction}{0.75}
\setlength{\parindent}{0pt}\setlength{\parskip}{0.55em}
\titleformat{\section}{\headfont\Large}{\thesection}{0.8em}{}
\titleformat{\subsection}{\headfont\large}{\thesubsection}{0.8em}{}
\captionsetup{font=small,labelfont=bf}
\urlstyle{same}
\hypersetup{pdftitle={Enclosure and the Swing Riots},pdfauthor={Tristan Hotham}}
\begin{document}
\begin{center}
{\LARGE\bfseries Enclosure and the Swing Riots\par}
\vspace{0.4em}
{\large A pre-specified parish-level test of the grievance hypothesis, England and Wales, 1830 to 1831\par}
\vspace{1.2em}
{\large Tristan Hotham\par}
\vspace{0.3em}
{\small Independent researcher\\ ORCID: 0000-0003-4931-9349\\ Data and code: https://doi.org/10.5281/zenodo.24091976\par}
\vspace{0.8em}
2026
\end{center}

\begin{quote}\small
\begin{center}\textbf{Abstract}\end{center}
Hobsbawm and Rud\'e argued that parliamentary enclosure helped to cause the Swing riots of 1830; Caprettini and Voth showed that threshing machines mattered. I link {{n.places_all|int}} places named in {{n.acts_all|int}} enclosure Acts to the parishes in Caprettini and Voth's data ({{n.matched_window_pc|f0}}\% of places named in Acts of 1790 to 1830 linked) and test, in a model specified in advance with county fixed effects and their controls, whether recent enclosure predicts unrest in {{n.sample|int}} parishes. As planned, parishes named in an Act of 1790 to 1830 were {{main.recent.b|f1}} percentage points more likely to see an incident (95\% CI {{main.recent.lo|f1}} to {{main.recent.hi|f1}}; mean {{main.recent.ymean|f1}}\%). Three findings argue against a grievance reading. Parishes enclosed only before 1790 show the same gap ({{main.old.b|f1}} points); stricter linking rules shrink it to between {{s.link_agree.b|f1}} and {{s.link_name.b|f1}} points; and in exploratory models that add parish area, a control the plan omitted, it disappears ({{s.x_area.b|f1}} points, 95\% CI {{s.x_area.lo|f1}} to {{s.x_area.hi|f1}}); larger parishes were both more often enclosed and more likely to record any incident. The planned estimate is therefore not robust evidence that recent parliamentary enclosure predicted where Swing broke out.

\textbf{Keywords:} enclosure; Swing riots; collective action; agrarian protest; parliamentary enclosure Acts; record linkage; pre-registration; England
\end{quote}

\section{Introduction}
Between late August 1830 and the spring of 1831, agricultural labourers in some {{n.riot_parishes|int}} parishes of the main sample studied here, most of them in southern and eastern England, broke threshing machines, burned stacks and barns, demanded higher wages and lower tithes, and sent threatening letters over the name of ``Captain Swing''. The rising was put down by special commissions that sentenced many of those tried to imprisonment, transportation or death \citep{hobsbawmrude1969}. Ever since the Hammonds' \textit{The Village Labourer}, written history has tied the misery behind Swing to the enclosure of open fields and commons in the previous half century \citep{hammond1911}. In the most influential modern account, Hobsbawm and Rud\'e placed enclosure among the long-run forces, with the poor law, low wages and seasonal unemployment, that turned the southern labourer into a wage-dependent proletariat ready to riot when the harvest of 1830 was in and the threshing machines were set to work \citep{hobsbawmrude1969}.

Quantitative work has since moved the emphasis. Using a newly digitised census of Swing incidents \citep{holland2005} and parish-level records of threshing machines from newspaper advertisements and agricultural surveys, \citet{caprettinivoth2020} showed that parishes where machines had spread were markedly more likely to see unrest, and they used the suitability of soils for threshing to argue that the relation is causal. \citet{aidtleon2022} showed that the rising spread along personal and trade networks. In this quantitative work enclosure has not been measured Act by Act; Caprettini and Voth's data carry only a district-level estimate drawn from \citet{gonner1912}.

This paper asks a narrower question that the new data make answerable. Using the Act-by-Act record of parliamentary enclosure, does a parish's own recent enclosure predict Swing unrest there, beyond threshing machines and the other covariates? The test requires linking the places named in enclosure Acts to the parishes of the 1851 census used by \citet{caprettinivoth2020}, which I do by name within county and by gazetteer coordinates, reporting the linkage rate and how results move under four alternative linking rules. Because the question has a long history and many plausible specifications, I wrote the analysis plan, including the hypotheses, the main specification and what would count against the hypothesis, before estimating any model of the outcome, and I report the planned specification as the main result.

The planned model finds a small positive association: parishes named in an enclosure Act of 1790 to 1830 were {{main.recent.b|f1}} percentage points more likely to see at least one Swing incident than parishes with no Act before 1831, against a mean of {{main.recent.ymean|f1}}\%. Three findings argue against reading this as support for the grievance hypothesis. First, there is no gradient in time: parishes whose only Acts predate 1790, two generations before the rising, show the same association, which contradicts the timing that a grievance mechanism implies. Second, the association is fragile to how Acts are linked to parishes. Third, it disappears when parish area, a variable that the planned model omitted, is held constant. Larger parishes were both more likely to be named in an Act and more likely to record at least one incident of any kind, so the planned estimate most plausibly reflects parish size. I label the area models exploratory because they were added after the planned results were seen.

The contribution is threefold. It provides the first parish-level linkage of the parliamentary enclosure Acts to the Swing data, with an audit and published link table. It offers a transparent, pre-specified test of a central claim of the social history of rural protest, with the robustness, placebo and leave-one-out checks expected of observational work. And it identifies a general hazard for parish-level studies of collective action: when the outcome is ``any incident in the parish'', parish size is a confounder that standard density controls do not remove.

\section{Background and literature}
\subsection{Parliamentary enclosure}
Enclosure extinguished common rights over open arable fields, meadows and wastes and replaced scattered strips with consolidated holdings in severalty. Much of England was enclosed piecemeal or by agreement before 1700, and in some southern and eastern counties little open field remained by the eighteenth century \citep{wordie1983}. From the 1750s, enclosure was increasingly carried out by private Act of Parliament, which allowed the owners of most of the land by value to override objectors \citep{turner1980,mingay1997}. Tate's \textit{Domesday} of enclosure Acts and awards is the standard county-by-county finding list \citep{tate1978}, and the enclosure maps that recorded the new allotments survive for thousands of parishes \citep{kain2004}. Parliamentary enclosure came in two waves, in the 1760s and 1770s, concentrated in the open-field Midlands, and during the French wars of 1793 to 1815, when high grain prices drew in heavier clays, wastes and commons across a wider area \citep{turner1980,turner1984}. Bogart and Richardson treat enclosure Acts as one instance of Parliament's role in reorganising property rights \citep{bogartrichardson2011}.

The economic effects remain debated. \citet{allen1982} argued that enclosure mainly redistributed rent from farmers and commoners to landowners with modest efficiency gains, and \citet{allen1992} that the yeoman's agricultural revolution in the South Midlands predated enclosure. Recent work using the Act data finds higher agricultural productivity and inequality in enclosed parishes \citep{heldring2022}. On labour, \citet{chambers1953} argued that enclosure increased rather than reduced the demand for rural labour. Against that, \citet{humphries1990} and \citet{neeson1993} showed how common rights supported labouring families, especially women and children, and how their loss reduced household independence; \citet{snell1985} found that seasonal unemployment and the decline of service deepened in the south after 1780. \citet{shawtaylor2001} argued that by the late eighteenth century most southern labourers had few common rights left to lose, so that parliamentary enclosure could not have created the agricultural proletariat, while \citet{clarkclark2001} estimated that common land was a small share of English land by the eighteenth century. This debate bears directly on the present test: if most labourers in the Swing counties had been landless for generations, a parish-specific shock from a recent Act should be small.

\subsection{The Swing riots}
\citet{hobsbawmrude1969} mapped the rising and catalogued its forms: machine breaking, arson, wage riots, tithe riots, threatening letters, robbery and assaults on overseers. They emphasised long-run proletarianisation, the allowance system of poor relief \citep{boyer1990}, and the immediate trigger of the 1830 harvest, the July Revolution in France and the general election. \citet{charlesworth1979} studied the spatial diffusion of the disturbances along roads and market networks, and his \textit{Atlas of Rural Protest} placed Swing in the longer history of rural unrest \citep{charlesworth1983,charlesworth1994}. Later work stressed continuity with incendiarism and other covert protest \citep{archer1990}, the political vocabulary of the rioters \citep{griffin2012}, and the wider landscape of popular contention \citep{tilly1995,stevenson1979,dunbabin1974}; \citet{randall2009} reviews the field forty years on. The broader frame comes from Thompson's accounts of plebeian custom and the moral economy \citep{thompson1963,thompson1971}.

The Family and Community Historical Research Society's project rebuilt the record of incidents parish by parish \citep{holland2005}, and this census underlies the quantitative work. \citet{caprettinivoth2020} combined it with threshing-machine records and 1801 census covariates for 1851 parishes, and found that a parish with a machine was considerably more likely to riot, using soil suitability as an instrument. \citet{aidtleon2022} modelled diffusion through networks, and \citet{aidtfranck2015} used the riots as a measure of the threat of revolution behind the Reform Act of 1832. Neither study models enclosure Act by Act. Caprettini and Voth's data include a district-level estimate of land enclosed by 1820 from \citet{gonner1912}, available for about two thirds of parishes.

\subsection{Measuring protest and exposure at the parish level}
Parish-level analysis of Swing has two strengths and two weaknesses. Its strength is that the parish was the unit of the poor law, of tithe and of enclosure, and the arena in which labourers, farmers and overseers met; many incidents were directed at named local targets \citep{hobsbawmrude1969,griffin2012}. It also allows comparisons within counties that hold constant the regional trajectory of the rising. Its weaknesses are that labour markets and protest networks crossed parish boundaries \citep{charlesworth1979,aidtleon2022}, so that parish exposure can mismeasure the exposure of the people who rioted, and that parishes differ greatly in size, so that the probability of recording an event depends on how much territory and how many people a parish holds. The first weakness biases parish estimates towards zero when the true effect works at district scale; the second can bias them in either direction, depending on how size relates to the exposure. The present study meets the first by testing spillovers to neighbouring parishes and the second, after the fact, by controlling for area.

Exposure to enclosure is also harder to measure than it seems. The Act names the parish or township whose fields were enclosed, but the land affected might be a small common or the whole open field, and the people affected might live in a neighbouring parish with rights of common. Tate's \textit{Domesday} and Turner's county tabulations give acreages for many Acts \citep{tate1978,turner1980}, but no machine-readable version linked to parishes was available for this study, so the exposure here is binary.

\subsection{What a grievance mechanism predicts}
If enclosure fed Swing through the grievance of lost common rights, three patterns should appear in parish data. Parishes enclosed shortly before 1830 should riot more than comparable unenclosed parishes; the association should fade with time since enclosure, as households adjusted and a new generation grew up without rights to remember; and it should be stronger for forms of protest aimed at farmers, landlords and overseers (arson, wage and relief riots, threatening letters) than for attacks on threshing machines. Spillovers are also possible, if labourers moved between neighbouring parishes or if enclosure in a district depressed wages there. Against these, enclosure may mark parishes that differed in other ways: open-field parishes were often large, nucleated, arable and populous, all features that might raise the chance of any incident.

\section{Data}
\subsection{Parishes and unrest}
The unit of analysis is the parish of the 1851 census as defined in Caprettini and Voth's replication data \citep{caprettinivoth2020data}, built on the 1851 parish boundaries of \citet{satchell2016}. The data hold {{n.parishes_all|int}} parishes of England and Wales; their main sample, which I use throughout, has {{n.sample|int}}. For each parish the data give the number of Swing incidents from \citet{holland2005}, split into attacks on threshing machines and other incidents; recorded threshing machines before 1830; population density, the agricultural share of occupations and the sex ratio in 1801; distances to Elham in Kent, where the rising began, and to the nearest town with a newspaper; cereal suitability and the share of heavy soils; and, for subsamples, Gonner's district enclosure estimate, poor rates in 1803 and replies to the 1834 Poor Law inquiry. In the main sample {{n.riot_parishes|int}} parishes ({{n.riot_pc|f1}}\%) saw at least one incident, {{n.incidents|int}} incidents in all; {{n.machine_attack_parishes|int}} saw machine breaking and {{n.other_parishes|int}} other forms of unrest. {{n.machines|int}} parishes had a recorded threshing machine.

\subsection{Enclosure Acts}
The Acts come from my earlier compilation of every enclosure Act listed in the Chronological Tables of Private and of Local Acts on legislation.gov.uk \citep{hotham2026enclosure}: {{n.acts_all|int}} Acts from {{n.first_act|raw}} to {{n.last_act|raw}}, naming {{n.places_all|int}} places, each with the year of the session and the county or counties named in the title, and, where it could be found, a match to a Wikidata place with coordinates. Acts peak in {{n.acts_peak_year|raw}}; {{n.acts_1793_1815|int}} fall in the war years 1793 to 1815. The tables list Acts, not acreage or awards, and do not cover enclosure by agreement or by the Inclosure Commissioners after 1845. Year is the first calendar year of the parliamentary session. The award that implemented an Act usually followed one to five years later \citep{turner1980}, so an Act of 1828 may not have been implemented by 1830; this affects only the shortest window.

\subsection{Linking Acts to parishes}
Each place named in an Act was linked to a parish by a rule fixed before any outcome model was run (Appendix~\ref{app:link}). In short: a place is linked by name when its name, an alternative spelling in the title, or its gazetteer name has the same normalised key as a parish (or as a component of a compound parish name such as ``Stow cum Quy'') in the county named in the title; failing that, simple variants are tried (``Woodend in Blakesley'' to Blakesley, ``Speen Moor'' to Speen). When several parishes share the key, the one nearest the gazetteer point is taken. If no name matches but the place has gazetteer coordinates, it is linked to the parish whose centroid is nearest, provided it lies within the larger of 2 km and 1.2 times that parish's equivalent radius. After reading cases by eye and before any outcome model, I added two rules: a name match more than 10 km from the gazetteer point is not used ({{n.conflict|int}} places, such as one of many Yorkshire Carltons), and a unique name match for a place the gazetteer found several of in its county and could not place is set aside ({{n.ambiguous|int}} places).

{{TABLE_LINK}}

Of the {{n.places_window|int}} places named in Acts of 1790 to 1830, {{n.matched_window|int}} ({{n.matched_window_pc|f1}}\%) are linked, {{n.name_window|int}} by name and {{n.spatial_window|int}} by point; {{n.acts_window_linked|int}} of the {{n.acts_window_named|int}} Acts of those years ({{n.acts_window_linked_pc|f1}}\%) have at least one linked parish, and {{n.in_sample_pc|f0}}\% of linked places fall in parishes of the main sample. Where a place has both a name link and a gazetteer point ({{n.name_point_both|int}} places in all years), the two agree on the parish in {{n.name_point_agree_pc|f1}}\% of cases, and the median distance from the gazetteer point to the centroid of the name-linked parish is {{n.name_point_median_km|f2}} km. A random audit of 100 links for 1790 to 1830, read against the titles and the parish list, found {{n.audit_ok|int}} correct or correct assignments of a township to its parish, and three doubtful. Unlinked places are mostly townships, hamlets and named commons absent from the parish list, and misspellings in the titles (Table~\ref{tab:link}). A parish is coded by the years of all Acts linked to it.

\subsection{Exposure}
The main exposure, as planned, is \textit{recent} enclosure: the parish is named in at least one Act of 1790 to 1830. \textit{Old-only} parishes are named only in Acts before 1790. The reference group is parishes not named in any Act before 1831. In the main sample, {{n.recent|int}} parishes ({{n.recent_pc|f1}}\%) are recent, {{n.old_only|int}} old-only and {{n.never|int}} never enclosed by Act before 1831; {{n.future|int}} of the latter were first named in an Act of 1831 to 1870, which provides a placebo. Table~\ref{tab:desc} compares the groups and Figure~\ref{fig:map} maps them.

{{TABLE_DESC}}

\begin{figure}[tbp]\centering\includegraphics[width=\textwidth]{f_map.png}
\caption{Left: parishes of the main sample by enclosure Act. Right: parishes with Swing incidents, 1830 to 1831, sized by number of incidents. Parish centroids; county outlines approximate.}\label{fig:map}\end{figure}

As a check on the linkage, the Act-based measure can be compared with the district estimate of land enclosed by 1820 that Caprettini and Voth take from \citet{gonner1912}. The two measure different things: Gonner's figure is a share of land, for a district, and includes non-parliamentary enclosure, whereas mine records whether a particular parish was named in an Act. At the parish level the correlation between Gonner's measure and being named in any Act before 1831 is weak ($r={{validation.r|f2}}$, {{validation.n|int}} parishes), as it must be when one variable is constant within districts; across the {{validation.districts|int}} market-town districts with Gonner data, the share of parishes named in an Act and Gonner's estimate are moderately correlated (Spearman $\rho={{validation.rho|f2}}$). The Act data therefore track the broad geography of enclosure that Gonner recorded while adding variation within districts, which is what a parish-level test needs.

The two maps already show the difficulty. Parliamentary enclosure was densest in the Midland belt and in Norfolk; Swing was densest in Kent, Sussex, Hampshire, Wiltshire, Berkshire and Norfolk. Many southern parishes where Swing struck hardest had been enclosed long before by agreement \citep{wordie1983}, so ``no Act'' there does not mean open fields survived. Comparisons across counties mix the two geographies; the design therefore compares parishes within counties.

\section{Research design}
The analysis plan was written in the study notes before any outcome model was run, after looking only at variable names, the sample size and the match-type counts of the gazetteer. It is reproduced in Appendix~\ref{app:plan}.

\textbf{Hypotheses.} H1: parishes enclosed by Act shortly before 1830 were more likely to see Swing unrest, holding threshing machines and the covariates of \citet{caprettinivoth2020} constant. H1b: the association is stronger for unrest other than machine breaking. H2: recent enclosure matters more than old enclosure.

\textbf{Main specification.} For parish $i$ in county $c$,
\[ \textit{Swing}_{ic} = \beta_1\,\textit{Recent}_{ic} + \beta_2\,\textit{OldOnly}_{ic} + X_{ic}'\gamma + \alpha_c + \varepsilon_{ic}, \]
where \textit{Swing} indicates at least one incident, $X$ holds log population density, agricultural share and log sex ratio in 1801, log distances to Elham and to the nearest newspaper town (the controls of Caprettini and Voth), an indicator for a recorded threshing machine, cereal suitability and the heavy-soil share, and $\alpha_c$ are {{n.counties|int}} county fixed effects (the three Yorkshire ridings separately). The model is a linear probability model estimated by OLS with standard errors clustered by county \citep{cameronmiller2015,abadie2023}. $\beta_1$ is the coefficient of interest, reported in percentage points with a 95\% confidence interval. The design is observational: $\beta_1$ is the difference in the chance of unrest between recently enclosed and unenclosed parishes of the same county with the same covariates, not the causal effect of enclosure.

\textbf{Planned robustness.} (1) A logit with the same controls and fixed effects (average marginal effects) and a Poisson model of the number of incidents. (2) Windows of 10, 20 and 40 years before 1830 and categories of years since the last Act. (3) Region or market-town fixed effects in place of county effects, and standard errors clustered on 50 km grid cells as a guard against spatial correlation across county lines \citep{conley1999}. (4) Placebos: enclosure first enacted in 1831 to 1870 among parishes with no earlier Act, and randomisation inference that permutes the exposure within counties 1,000 times \citep{young2019}. (5) Leave-one-county-out and leave-one-region-out estimates. (6) Linking sensitivity: name links only, gazetteer points only, links where name and point agree, an extended rule that adds close spellings, counties where at least 80\% of places were linked, and a simulation that moves 10\% or 25\% of linked places to a random neighbouring parish within 5 km. (7) Adding Gonner's district measure and, for the subsample with it, poor relief per head in 1803. (8) Neighbouring enclosure within 5 and 10 km.

\textbf{What would count against H1.} The plan stated that H1 would be contradicted if $\beta_1$ were not significantly positive and its confidence interval excluded effects of about 3 points (a fifth of the mean); if the future-enclosure placebo showed an association of similar size; or if the estimate were driven by a single county.

\section{Results}
\subsection{The planned estimate}
Table~\ref{tab:main} reports the main results. Without controls or fixed effects, recent enclosure is associated with a {{s.raw.b|f1}} point higher chance of unrest. Within counties and with the planned controls, the coefficient is {{main.recent.b|f2}} points (95\% CI {{main.recent.lo|f2}} to {{main.recent.hi|f2}}; $p={{main.recent.p|f3}}$; $N={{main.recent.n|int}}$), about {{main.rel_pc|f0}}\% of the mean of {{main.recent.ymean|f1}}\%. By the criteria set in advance, H1 is not rejected by the main specification.

{{TABLE_MAIN}}

The control coefficients behave as expected (Appendix Table~\ref{tab:full}). A recorded threshing machine is associated with a {{full.main.coefs.thresh1.b|f3}} higher probability of unrest, close in size to the descriptive gap in \citet{caprettinivoth2020}; denser parishes and those with lower agricultural shares were more likely to see incidents.

\textbf{Precision.} The cluster-robust standard error of the main coefficient is {{T.se_main|f2}} points, so the minimum effect detectable with 80\% power at the 5\% level is about {{T.mde|f1}} points (2.8 standard errors). The 3-point threshold written into the plan is close to this, which means the design is informative about effects of the size implied by accounts in which enclosure is a principal cause of the rising, but not about effects of one or two points. This was not computed before the analysis and is reported to help readers judge the null results below.

\subsection{Descriptive patterns}
Before turning to the timing, it is worth reading Table~\ref{tab:desc} closely. In raw terms, {{desc.recent.riot|f1}}\% of recently enclosed parishes saw unrest, against {{desc.old_only.riot|f1}}\% of old-only and {{desc.never.riot|f1}}\% of never-enclosed parishes, a raw gap of {{desc.gap_raw|f1}} points between recent and never. The recent group also had more machine attacks ({{desc.recent.machine_attack|f1}}\% against {{desc.never.machine_attack|f1}}\%), although threshing machines were recorded in about the same share of parishes in each group ({{desc.recent.machines|f1}}\% and {{desc.never.machines|f1}}\%). The similarity in machine records matters: if parliamentary enclosure had encouraged larger arable farms that adopted threshing machines, enclosure would act partly through machines and controlling for machines would understate its total association. The recorded machines give no sign of such a channel, though the records are a sample. The {{n.future|int}} parishes enclosed only after 1830 saw unrest at {{desc.future.riot|f1}}\%, slightly below the never-enclosed group.

Within counties the picture is mixed. In Norfolk, where the second wave of Acts coincided with heavy Swing activity, enclosed and unenclosed parishes both rioted often; in Kent and Sussex, the heartland of the rising, relatively few parishes were named in an Act of 1790 to 1830 ({{cty.KENT.recent_pc|f0}}\% and {{cty.SUSSEX.recent_pc|f0}}\%, against {{cty.NORFOLK.recent_pc|f0}}\% in Norfolk); in the Midland counties with the most Acts, Swing was thin. The county fixed effects therefore matter a great deal: they remove the negative cross-county association between the density of Acts and the density of riots and leave the within-county comparison that the hypothesis is about. The raw coefficient (column 1 of Table~\ref{tab:main}) and the model with controls but without county effects ({{s.ctrl_nofe.b|f1}} points, Table~\ref{tab:rob}) are larger than those with county effects, which shows that some of the raw gap comes from the counties in which enclosed parishes lie.

\subsection{Timing}
H2 fails. Old-only parishes show a coefficient of {{main.old.b|f2}} points (95\% CI {{main.old.lo|f2}} to {{main.old.hi|f2}}), as large as for recent enclosure; the difference, recent minus old, is {{main.h2_diff|f2}} points (95\% CI {{main.h2_lo|f2}} to {{main.h2_hi|f2}}; $p={{main.h2_p|f2}}$). Treated as one category, any Act before 1831 carries {{anyact.b|f2}} points. The window variants tell the same story (Table~\ref{tab:timing}). Acts within 10 years of 1830 carry the largest point estimate, {{windows.w10.win.b|f1}} points, but it is imprecise (95\% CI {{windows.w10.win.lo|f1}} to {{windows.w10.win.hi|f1}}), and parishes whose Acts were all earlier than the window carry about the same coefficient as the window itself at 20 and 40 years. By years since the last Act, the categories 20 to 39 and 40 or more years before 1830 are about as large as the most recent ones. Figure~\ref{fig:timing} shows the raw shares by decade of the last Act: parishes enclosed in the 1820s rioted most often ({{decade.riot.7|f1}}\%, $n={{decade.n.7|int}}$), but parishes last enclosed in the 1780s rioted as often as those of the 1800s.

{{TABLE_TIMING}}

\begin{figure}[tbp]\centering\includegraphics[width=0.92\textwidth]{f_timing.png}
\caption{Share of parishes with any Swing incident by decade of their last enclosure Act before 1831 (decades with at least 30 parishes), against parishes with no Act. Raw shares without controls.}\label{fig:timing}\end{figure}

\subsection{Forms of unrest}
H1b receives no clear support. For machine breaking, the recent coefficient is {{s.machine.b|f2}} points (95\% CI {{s.machine.lo|f2}} to {{s.machine.hi|f2}}) against a mean of {{s.machine.ymean|f1}}\%; for other unrest it is {{s.other.b|f2}} points (95\% CI {{s.other.lo|f2}} to {{s.other.hi|f2}}) against a mean of {{s.other.ymean|f1}}\%. Relative to their means the two are similar, and old-only enclosure is again as large for other unrest ({{s.other_old.b|f2}} points). In the Poisson model for the number of incidents the incidence rate ratio is {{s.poisson.irr|f2}} (95\% CI {{s.poisson.irr_lo|f2}} to {{s.poisson.irr_hi|f2}}) for recent and {{s.poisson_old.irr|f2}} for old-only enclosure.

\subsection{Placebo, inference and influence}
Parishes first named in an Act of 1831 to 1870, which cannot have been affected by enclosure before the riots, show a coefficient of {{placebo.b|f1}} points (95\% CI {{placebo.lo|f1}} to {{placebo.hi|f1}}; $N={{placebo.n|int}}$ parishes, {{n.future|int}} of them future-enclosed). The placebo is imprecise, but it is not of similar size and has the opposite sign, so the planned criterion is met. Permuting the exposure within counties gives a randomisation $p$ of {{ri.p|f3}} for $\beta_1$ (the central 95\% of permuted coefficients lies between {{ri.lo|f2}} and {{ri.hi|f2}}), smaller than the cluster-robust $p$, as expected given that permutation within counties ignores spatial correlation. With standard errors clustered on 50 km cells the 95\% interval is {{s.grid_se.lo|f2}} to {{s.grid_se.hi|f2}}.

No single county drives the estimate. Leaving out each of the {{loo.n|int}} counties in turn gives coefficients from {{loo.min|f2}} (without {{loo.min_county|raw}}) to {{loo.max|f2}} (without {{loo.max_county|raw}}); all are positive and {{loo.n_sig|int}} have $p<0.05$ (Figure~\ref{fig:loo}). Leaving out regions gives {{lor_range.min|f2}} to {{lor_range.max|f2}} (Appendix Table~\ref{tab:lor}).

\begin{figure}[tbp]\centering\includegraphics[width=0.92\textwidth]{f_loo.png}
\caption{The main coefficient on recent enclosure re-estimated leaving out one county at a time, with 95\% confidence intervals.}\label{fig:loo}\end{figure}

\section{Robustness}
\subsection{Specification}
Figure~\ref{fig:coef} and Table~\ref{tab:rob} collect every specification. Logit marginal effects ({{s.logit.b|f2}} points), market-town fixed effects ({{s.market_fe.b|f2}}) and region fixed effects ({{s.region_fe.b|f2}}) are close to the main estimate, and controlling for neighbouring enclosure within 5 and 10 km leaves it at {{s.spill.b|f2}}; the neighbour indicators themselves are small and negative ({{s.spill5.b|f2}} and {{s.spill10.b|f2}} points), which gives no sign of district-wide spillovers. Controlling for pre-1830 unrest changes nothing. In the two-thirds of parishes covered by Gonner's district estimates the coefficient falls to {{s.gonner_base.b|f2}} (95\% CI {{s.gonner_base.lo|f2}} to {{s.gonner_base.hi|f2}}) before Gonner's measure is added and {{s.gonner.b|f2}} after; in the {{s.poor_base.n|int}} parishes with poor rates for 1803 it is {{s.poor_base.b|f2}}, with a wide interval.

\begin{figure}[tbp]\centering\includegraphics[width=0.95\textwidth]{f_coef.png}
\caption{Coefficient on recent enclosure (Act 1790 to 1830) across specifications, percentage points with 95\% confidence intervals. Black: planned specifications. Blue: alternative linking rules. Red: exploratory checks added after the planned results were seen.}\label{fig:coef}\end{figure}

\subsection{Linking}
The estimate is sensitive to how Acts are linked to parishes. Using only name links gives {{s.link_name.b|f2}} points (95\% CI {{s.link_name.lo|f2}} to {{s.link_name.hi|f2}}); only gazetteer points, {{s.link_spatial.b|f2}}; only links where name and point agree, {{s.link_agree.b|f2}} (95\% CI {{s.link_agree.lo|f2}} to {{s.link_agree.hi|f2}}). The extended rule with close spellings gives {{s.link_ext.b|f2}}, and the {{n.hi_counties|int}} counties where at least 80\% of places were linked give {{s.hi_match.b|f2}} (95\% CI {{s.hi_match.lo|f2}} to {{s.hi_match.hi|f2}}). Each strict rule moves some truly enclosed parishes into the reference group, which attenuates the estimate, so smaller coefficients under strict rules are expected; but none is significant, and the simulation shows that random linking errors of the kind the main rule may make would shift the estimate only modestly (mean {{sim.s10.mean|f2}} and {{sim.s25.mean|f2}} points when 10\% and 25\% of places are moved to a neighbouring parish). Dropping the {{s.x_clean_dropped|int}} parishes whose exposure depends on the linking rule gives {{s.x_clean.b|f2}} points (95\% CI {{s.x_clean.lo|f2}} to {{s.x_clean.hi|f2}}).

\subsection{Parish size (exploratory)}
One omission from the plan turns out to matter. The controls of \citet{caprettinivoth2020}, which the plan adopted, include population density but no measure of parish area or of total population. The outcome is whether any incident was recorded in the parish, and the 95th percentile of parish area is about {{n.area_ratio_95_5|f0}} times the 5th. Larger parishes have more farms, more ricks and more people to be counted, so the chance of at least one incident rises with size: from {{area_q.top_bottom_riot.0|f1}}\% in the smallest fifth of parishes by area to {{area_q.top_bottom_riot.1|f1}}\% in the largest (Figure~\ref{fig:size}). Larger parishes were also more often named in an enclosure Act, from {{area_q.enc_share.0|f0}}\% to {{area_q.enc_share.4|f0}}\%; the median area is {{area.recent|f1}} km\textsuperscript{2} for recent, {{area.old_only|f1}} km\textsuperscript{2} for old-only and {{area.never|f1}} km\textsuperscript{2} for never-enclosed parishes. Population density, which the planned model includes, does not absorb this, because a large parish can be thinly peopled and still hold more people, and more targets, than a small one.

\begin{figure}[tbp]\centering\includegraphics[width=0.95\textwidth]{f_size.png}
\caption{Parish area, Swing unrest and enclosure. Fifths of the main sample by area of the 1851 parish.}\label{fig:size}\end{figure}

Adding log parish area to the main model reduces the recent coefficient to {{s.x_area.b|f2}} points (95\% CI {{s.x_area.lo|f2}} to {{s.x_area.hi|f2}}) and the old-only coefficient to {{s.x_area_old.b|f2}}; log area itself carries {{s.x_area_logarea.b|f1}} points per log unit. Area deciles give {{s.x_area_dec.b|f2}} (95\% CI {{s.x_area_dec.lo|f2}} to {{s.x_area_dec.hi|f2}}), a logit with log area {{s.x_area_logit.b|f2}}, and a Poisson model of the count an incidence rate ratio of {{s.x_area_poisson.irr|f2}}. For unrest other than machine breaking, the form most directly tied to grievances against farmers and overseers, the coefficient with area is {{s.x_area_other.b|f2}}. Even the 10-year window falls to {{windows.w10_area.b|f1}} points (95\% CI {{windows.w10_area.lo|f1}} to {{windows.w10_area.hi|f1}}). These intervals exclude the 3-point effect named in the plan as the threshold of substantive interest, except for the imprecise 10-year window. The size control also moves the future-enclosure placebo: with log area, parishes first named in an Act of 1831 to 1870 were {{placebo_area.b|f1}} points less likely to see unrest than never-enclosed parishes of the same size (95\% CI {{placebo_area.lo|f1}} to {{placebo_area.hi|f1}}; $p={{placebo_area.p|f3}}$). I have no ready explanation for this; it shows that parishes named in Acts at different dates differ in ways the controls do not capture, and it is a further reason for caution about any parish-level enclosure contrast. Because these models were added after seeing the planned results, they should be read as exploratory, but the mechanism is simple and the change is large.

\subsection{Heterogeneity (exploratory)}
The association does not differ by whether the parish had a recorded threshing machine (interaction {{s.het_machine.b|f2}} points, 95\% CI {{s.het_machine.lo|f2}} to {{s.het_machine.hi|f2}}) or between the South-East region and the rest ({{s.het_se.b|f2}}, 95\% CI {{s.het_se.lo|f2}} to {{s.het_se.hi|f2}}).

\section{Discussion}
The planned test finds that parishes enclosed by Act in the forty years before Swing were somewhat more likely to see unrest, which is what a reader of Hobsbawm and Rud\'e might expect. The rest of the evidence makes that reading hard to sustain. A grievance born of the loss of common rights should be sharper where the loss was recent. It is not: parishes enclosed only before 1790, whose labourers in 1830 would mostly have been born after enclosure, show the same association. The estimate weakens when the linkage is made stricter, and it vanishes once parish size, which is associated with both enclosure and the chance of recording any incident, is held constant.

These results fit the argument of \citet{shawtaylor2001} that most southern labourers had lost access to land well before the parliamentary enclosure of their parish, and the view that the immediate drivers of 1830 were wages, winter unemployment and the poor law \citep{snell1985,boyer1990}, with threshing machines as a specific threat to winter work \citep{caprettinivoth2020}. They do not show that enclosure was irrelevant. Several mechanisms in the literature operate at scales the parish test cannot see. Enclosure may have lowered wages or raised poor rates across a labour market rather than a parish, and labourers moved between neighbouring parishes; the neighbour indicators, which are negative, speak against a simple spillover, but districts may be larger than 10 km. Much enclosure in the Swing heartland was not parliamentary at all \citep{wordie1983}, so the comparison is between parishes enclosed by Act and parishes enclosed earlier by other means, not between enclosed and open parishes; in that sense the test is of the specific claim that the parliamentary enclosures of 1790 to 1830 left local scars, not of enclosure as a long process. And Swing incidents were not all equivalent: the outcome gives a large village with one arson the same weight as a small one where the whole parish turned out.

The finding on parish size has a methodological lesson beyond this case. Studies that regress ``any event in the unit'' on unit characteristics should control for unit size, or model counts with exposure offsets, since otherwise any characteristic correlated with size will appear to predict events. Density controls do not do this job. The size of the 1851 parish also enters the threshing-machine result, since machines were recorded in parishes and larger parishes had more farms, and in the model with log area the machine coefficient falls from {{full.main.coefs.thresh1.b|f3}} to {{full.x_area.coefs.thresh1.b|f3}}, though it remains significant ($p={{full.x_area.coefs.thresh1.p|f3}}$; Appendix Table~\ref{tab:full}). This bears only on the descriptive association; Caprettini and Voth's instrumental-variable estimate rests on a different argument, and I do not re-examine it here.

\subsection{How large would an enclosure effect have had to be?}
It helps to translate the estimates into incidents. Taken at face value, the planned coefficient implies that about {{T.implied_parishes|f0}} of the {{n.riot_parishes|int}} parishes with unrest in the main sample (roughly {{T.implied_share|f0}}\%) owed their unrest to recent parliamentary enclosure. The upper end of the 95\% interval in the model with parish area implies about {{T.implied_parishes_area_hi|f0}}. Either way, recent enclosure by Act could account for only a small part of the geography of Swing, which was dominated by the south-eastern counties where parliamentary enclosure was least common (Figure~\ref{fig:map}). By comparison, a recorded threshing machine is associated with a difference of {{full.main.coefs.thresh1.b|f3}} in probability in the same model, and the gradient with distance from Elham is steep \citep{caprettinivoth2020}. Whatever the role of enclosure in the making of the southern rural poor, the place-to-place pattern of the rising is not well predicted by where Parliament had recently enclosed.

\subsection{Why might old enclosure look like recent enclosure?}
Two explanations fit the equal coefficients for old and recent enclosure. The first, favoured by the size results, is that being named in an Act at any time marks a kind of parish: larger, often nucleated, with substantial arable and a body of proprietors able to promote a bill, and therefore with more potential targets and participants. The second is that enclosure set off slow changes in farm size, employment and settlement that persisted for generations, so that old enclosure mattered as much as new. The second explanation predicts that the old-enclosure coefficient should survive the size controls; it does not ({{s.x_area_old.b|f2}} points), which favours the first. Neither explanation supports the specific claim that the recent loss of common rights drove parishes into revolt in 1830.

\subsection{Relation to earlier quantitative work}
The descriptive analysis of the same parish data in an earlier study found that unrest varied little across Gonner's district enclosure fifths \citep{hotham2026swing}. The present paper confirms this at the parish level with Act data, adds the timing contrasts and placebo, and shows that a naive parish-level association would have pointed the other way. It also complements \citet{heldring2022}, who find lasting economic effects of parliamentary enclosure on productivity and inequality: effects on the economy of enclosed parishes need not translate into a higher propensity to riot in a particular year, and the riots of 1830 were concentrated in the south and east where the Acts studied by Heldring and co-authors were less dense.

\subsection{Limitations}
The design is observational. Parishes were not enclosed at random: enclosure required a majority of landowners by value and favoured particular soils, field systems and proprietors \citep{turner1980,allen1992}, and the county fixed effects and covariates cannot remove all differences between enclosed and other parishes. The Act data cover parliamentary enclosure only and record Acts, not the share of the parish enclosed or the acreage of common lost; an Act that enclosed a small common counts the same as one that enclosed the whole open field. Linkage reaches {{n.matched_window_pc|f0}}\% of places named in Acts of 1790 to 1830, and the unlinked places are not random: they are disproportionately townships and hamlets within large northern parishes and places with unusual spellings. The audit suggests that the links made are mostly right, but missing links place some enclosed parishes in the reference group, biasing coefficients toward zero; the simulation and the strict-link variants suggest this bias is modest relative to the size effect. The Swing record of \citet{holland2005} is the most complete available but depends on surviving newspaper and judicial reports, and the threshing-machine records are a sample of the machines in use \citep{caprettinivoth2020}. The year of an Act is not the year of the award, and the shortest window may include enclosures not yet carried out. Finally, the parish-size models were not in the plan; I report them because they change the interpretation, and readers should weigh them as exploratory.

\section{Conclusion}
Parliamentary enclosure has been part of the story of Captain Swing since the Hammonds. Linking the enclosure Acts to the parish data of the Swing literature allows a direct test. A pre-specified model finds that parishes named in an Act of 1790 to 1830 were {{main.recent.b|f1}} percentage points more likely to see unrest, but parishes enclosed long before show the same difference, stricter linking weakens it, and, in exploratory models added after the planned results were seen, holding parish area constant removes it. The parish-level data therefore give no support to the claim that recent parliamentary enclosure predicted where Swing broke out. Enclosure's role in the making of the rural proletariat, which unfolded over generations and across districts, is a different question that these data cannot settle.

\section*{Data availability statement}
Data and code: Zenodo DOI https://doi.org/10.5281/zenodo.24091976. The linked parish data (\texttt{parishes\_enclosure.csv}), the Act-to-parish link table with the method used for each place (\texttt{act\_parish\_links.csv}), the leave-one-county-out estimates and \texttt{results.json}, which holds every number in this paper, are included, with the code that rebuilds them from the replication data of \citet{caprettinivoth2020data} (CC BY 4.0) and the enclosure Act compilation \citep{hotham2026enclosure}. An interactive map is at \url{{{PAGE}}}.

\section*{AI statement}
{{AI}}

\begin{thebibliography}{99}\small
\bibitem[Abadie et al.(2023)]{abadie2023} Abadie, A., Athey, S., Imbens, G. W. and Wooldridge, J. M. (2023) When should you adjust standard errors for clustering? \textit{Quarterly Journal of Economics} 138(1), 1 to 35.
\bibitem[Aidt and Franck(2015)]{aidtfranck2015} Aidt, T. S. and Franck, R. (2015) Democratization under the threat of revolution: evidence from the Great Reform Act of 1832. \textit{Econometrica} 83(2), 505 to 547.
\bibitem[Aidt et al.(2022)]{aidtleon2022} Aidt, T. S., Leon, G. and Satchell, M. (2022) The social dynamics of collective action: evidence from the diffusion of the Swing riots, 1830 to 1831. \textit{Journal of Politics} 84(1), 209 to 225. doi:10.1086/714784.
\bibitem[Allen(1982)]{allen1982} Allen, R. C. (1982) The efficiency and distributional consequences of eighteenth century enclosures. \textit{Economic Journal} 92(368), 937 to 953.
\bibitem[Allen(1992)]{allen1992} Allen, R. C. (1992) \textit{Enclosure and the Yeoman: The Agricultural Development of the South Midlands, 1450 to 1850}. Oxford: Clarendon Press.
\bibitem[Archer(1990)]{archer1990} Archer, J. E. (1990) \textit{By a Flash and a Scare: Incendiarism, Animal Maiming and Poaching in East Anglia, 1815 to 1870}. Oxford: Clarendon Press.
\bibitem[Bogart and Richardson(2011)]{bogartrichardson2011} Bogart, D. and Richardson, G. (2011) Property rights and Parliament in industrializing Britain. \textit{Journal of Law and Economics} 54(2), 241 to 274. doi:10.1086/652901.
\bibitem[Boyer(1990)]{boyer1990} Boyer, G. R. (1990) \textit{An Economic History of the English Poor Law, 1750 to 1850}. Cambridge: Cambridge University Press.
\bibitem[Cameron and Miller(2015)]{cameronmiller2015} Cameron, A. C. and Miller, D. L. (2015) A practitioner's guide to cluster-robust inference. \textit{Journal of Human Resources} 50(2), 317 to 372.
\bibitem[Caprettini and Voth(2020a)]{caprettinivoth2020} Caprettini, B. and Voth, H.-J. (2020) Rage against the machines: labor-saving technology and unrest in industrializing England. \textit{American Economic Review: Insights} 2(3), 305 to 320. doi:10.1257/aeri.20190385.
\bibitem[Caprettini and Voth(2020b)]{caprettinivoth2020data} Caprettini, B. and Voth, H.-J. (2020) Data and code for: Rage against the machines. Ann Arbor: Inter-university Consortium for Political and Social Research, openICPSR project 117044, version V1. doi:10.3886/E117044V1.
\bibitem[Chambers(1953)]{chambers1953} Chambers, J. D. (1953) Enclosure and labour supply in the industrial revolution. \textit{Economic History Review} 5(3), 319 to 343.
\bibitem[Charlesworth(1979)]{charlesworth1979} Charlesworth, A. (1979) \textit{Social Protest in a Rural Society: The Spatial Diffusion of the Captain Swing Disturbances of 1830 to 1831}. Historical Geography Research Series 1. Norwich: Geo Abstracts.
\bibitem[Charlesworth(1983)]{charlesworth1983} Charlesworth, A., ed. (1983) \textit{An Atlas of Rural Protest in Britain, 1548 to 1900}. London: Croom Helm.
\bibitem[Charlesworth(1994)]{charlesworth1994} Charlesworth, A. (1994) The spatial diffusion of riots: popular disturbances in England and Wales, 1750 to 1850. \textit{Rural History} 5(1), 1 to 22.
\bibitem[Clark and Clark(2001)]{clarkclark2001} Clark, G. and Clark, A. (2001) Common rights to land in England, 1475 to 1839. \textit{Journal of Economic History} 61(4), 1009 to 1036.
\bibitem[Conley(1999)]{conley1999} Conley, T. G. (1999) GMM estimation with cross sectional dependence. \textit{Journal of Econometrics} 92(1), 1 to 45.
\bibitem[Dunbabin(1974)]{dunbabin1974} Dunbabin, J. P. D. (1974) \textit{Rural Discontent in Nineteenth-Century Britain}. London: Faber and Faber.
\bibitem[Gonner(1912)]{gonner1912} Gonner, E. C. K. (1912) \textit{Common Land and Inclosure}. London: Macmillan.
\bibitem[Griffin(2012)]{griffin2012} Griffin, C. J. (2012) \textit{The Rural War: Captain Swing and the Politics of Protest}. Manchester: Manchester University Press.
\bibitem[Hammond and Hammond(1911)]{hammond1911} Hammond, J. L. and Hammond, B. (1911) \textit{The Village Labourer, 1760 to 1832}. London: Longmans, Green.
\bibitem[Heldring et al.(2022)]{heldring2022} Heldring, L., Robinson, J. A. and Vollmer, S. (2022) The economic effects of the English parliamentary enclosures. NBER Working Paper 29772. Cambridge, MA: National Bureau of Economic Research.
\bibitem[Hobsbawm and Rud\'e(1969)]{hobsbawmrude1969} Hobsbawm, E. J. and Rud\'e, G. (1969) \textit{Captain Swing}. London: Lawrence and Wishart.
\bibitem[Holland(2005)]{holland2005} Holland, M., ed. (2005) \textit{Swing Unmasked: The Agricultural Riots of 1830 to 1832 and their Wider Implications}. Milton Keynes: FACHRS Publications.
\bibitem[Hotham(2026a)]{hotham2026enclosure} Hotham, T. (2026) Fencing the commons: mapping parliamentary enclosure in England and Wales from the Chronological Tables of Acts. Zenodo. doi:10.5281/zenodo.23873003.
\bibitem[Hotham(2026b)]{hotham2026swing} Hotham, T. (2026) Captain Swing: did the riots follow the machines? A parish-level look at the Swing riots of 1830 to 1832. Zenodo. doi:10.5281/zenodo.23727133.
\bibitem[Humphries(1990)]{humphries1990} Humphries, J. (1990) Enclosures, common rights, and women: the proletarianization of families in the late eighteenth and early nineteenth centuries. \textit{Journal of Economic History} 50(1), 17 to 42.
\bibitem[Kain et al.(2004)]{kain2004} Kain, R. J. P., Chapman, J. and Oliver, R. R. (2004) \textit{The Enclosure Maps of England and Wales, 1595 to 1918}. Cambridge: Cambridge University Press.
\bibitem[Mingay(1997)]{mingay1997} Mingay, G. E. (1997) \textit{Parliamentary Enclosure in England: An Introduction to its Causes, Incidence and Impact, 1750 to 1850}. London: Longman.
\bibitem[Neeson(1993)]{neeson1993} Neeson, J. M. (1993) \textit{Commoners: Common Right, Enclosure and Social Change in England, 1700 to 1820}. Cambridge: Cambridge University Press.
\bibitem[Randall(2009)]{randall2009} Randall, A. (2009) Captain Swing: a retrospect. \textit{International Review of Social History} 54(3), 419 to 427. doi:10.1017/S0020859009990393.
\bibitem[Satchell et al.(2016)]{satchell2016} Satchell, A. E. M., Kitson, P. K., Newton, G. H., Shaw-Taylor, L. and Wrigley, E. A. (2016) 1851 England and Wales census parishes, townships and places. UK Data Service ReShare.
\bibitem[Shaw-Taylor(2001)]{shawtaylor2001} Shaw-Taylor, L. (2001) Parliamentary enclosure and the emergence of an English agricultural proletariat. \textit{Journal of Economic History} 61(3), 640 to 662.
\bibitem[Snell(1985)]{snell1985} Snell, K. D. M. (1985) \textit{Annals of the Labouring Poor: Social Change and Agrarian England, 1660 to 1900}. Cambridge: Cambridge University Press.
\bibitem[Stevenson(1979)]{stevenson1979} Stevenson, J. (1979) \textit{Popular Disturbances in England, 1700 to 1870}. London: Longman.
\bibitem[Tate(1978)]{tate1978} Tate, W. E. (1978) \textit{A Domesday of English Enclosure Acts and Awards}, ed. M. E. Turner. Reading: University of Reading Library.
\bibitem[Thompson(1963)]{thompson1963} Thompson, E. P. (1963) \textit{The Making of the English Working Class}. London: Victor Gollancz.
\bibitem[Thompson(1971)]{thompson1971} Thompson, E. P. (1971) The moral economy of the English crowd in the eighteenth century. \textit{Past and Present} 50, 76 to 136.
\bibitem[Tilly(1995)]{tilly1995} Tilly, C. (1995) \textit{Popular Contention in Great Britain, 1758 to 1834}. Cambridge, MA: Harvard University Press.
\bibitem[Turner(1980)]{turner1980} Turner, M. E. (1980) \textit{English Parliamentary Enclosure: Its Historical Geography and Economic History}. Folkestone: Dawson.
\bibitem[Turner(1984)]{turner1984} Turner, M. E. (1984) \textit{Enclosures in Britain, 1750 to 1830}. London: Macmillan.
\bibitem[Wordie(1983)]{wordie1983} Wordie, J. R. (1983) The chronology of English enclosure, 1500 to 1914. \textit{Economic History Review} 36(4), 483 to 505.
\bibitem[Young(2019)]{young2019} Young, A. (2019) Channeling Fisher: randomization tests and the statistical insignificance of seemingly significant experimental results. \textit{Quarterly Journal of Economics} 134(2), 557 to 598.
\end{thebibliography}

\appendix
\section{Analysis plan}\label{app:plan}
The plan below was written in the study notes before any model of the outcome was run, and is reproduced with only typographical changes. Amendments to the linking rule made after reading link cases (but before any outcome model) are listed in Appendix~\ref{app:link}. Analyses not listed here are labelled exploratory in the text.

\begin{small}
\textit{Hypotheses.} H1 (Hobsbawm and Rud\'e): parishes enclosed by parliamentary Act shortly before 1830 were more likely to see Swing unrest in 1830 to 1831, holding constant threshing-machine exposure and the other covariates of Caprettini and Voth. H1b (mechanism, secondary): the association should be stronger for unrest other than machine breaking. H2 (timing): recent enclosure (Acts 1790 to 1830) should matter more than old enclosure (Acts before 1790).

\textit{Unit, sample, outcome.} The 1851 parish of Caprettini and Voth; their main sample (N = 9,674). Main outcome: any Swing incident (0/1). Secondary: the number of incidents; machine attacks and other incidents, each as any-incident indicators.

\textit{Exposure.} Link each place named in an enclosure Act to a parish: (1) normalised name equal to the parish name or a component of a compound name, within the county named in the title or the gazetteer county; (2) otherwise, if located, the parish whose centroid is nearest within max(2 km, 1.2 times the parish's equivalent radius). Main exposure: named in at least one Act of 1790 to 1830. Categories: recent, old only (Acts before 1790 only), never enclosed by Act before 1831 (reference). Windows of 10, 20 and 40 years. Spillover: other recent-Act parishes within 5 and 10 km.

\textit{Main specification.} Linear probability model of any incident on recent and old-only enclosure with the controls of Caprettini and Voth (log density, agricultural share, log sex ratio, log distance to Elham, log distance to newspaper town), threshing machine, cereal suitability and heavy-soil share, and county fixed effects; standard errors clustered by county, also on 50 km grid cells.

\textit{Robustness.} Logit (average marginal effects) and Poisson; windows and years since the last Act; region or market-town fixed effects; placebo with first Acts 1831 to 1870 among parishes not enclosed before 1831, and randomisation inference permuting exposure within county 1,000 times; leave-one-county-out; linking variants (name only, point only, agreement, counties with at least 80\% linked, simulated relinking of 10\% and 25\% of places); Gonner's district enclosure and 1803 poor relief in subsamples; heterogeneity by machine presence and by region (South-East against the rest).

\textit{What would count against H1.} The coefficient on recent enclosure in the main model is not significantly positive and its 95\% confidence interval excludes effects of about 3 percentage points; or the future-enclosure placebo shows an association of similar size; or the estimate is driven by one county.
\end{small}

\section{Linking details}\label{app:link}
Names on both sides are reduced to keys: accents, punctuation and bracketed alternatives (Welsh forms in the parish list) dropped, lower case, Saint to St, upon and super to on, juxta and near to next, magna to great, parva to little, cum to with, hyphens as spaces, ``the'' dropped. Counties named in titles are mapped to the parish data's counties, with Yorkshire spread over its three ridings and Hampshire including the Isle of Wight. A place with no county in its title and no gazetteer county is matched only if its key is unique in England and Wales. Variants strip land words from the end or start of a name (common, moor, marsh, fen, heath, field, waste, meadow, pasture, green, forest, chase, lordship, township, manor, hamlet and similar) and take the parish in ``X in Y''. The extended variant adds close spellings within the stated county (string similarity at least 0.85 with a clear winner); reading 25 of these showed both rescues (Ulverstone to Ulverston, Beaminister to Beaminster) and errors (Horning to Honing), so it is used only as a sensitivity check. The audit sample of 100 links and every link, with its method, are in \texttt{act\_parish\_links.csv}.

\section{Variables and data construction}\label{app:vars}
\textbf{Outcome.} \texttt{SWING} in the replication data counts incidents attributed to the parish in \citet{holland2005} for 1830 to 1832; \texttt{BREAKTHRESH} counts attacks on threshing machines and \texttt{REST} the other incidents. The main outcome is $\textit{Swing}=1$ if \texttt{SWING}$>0$.

\textbf{Controls.} Log population density in 1801 (persons per km\textsuperscript{2}); agricultural share of employment in 1801; log of the ratio of males to females in 1801; log distance to Elham (km); log distance to the nearest town with a newspaper (km); a threshing machine recorded in the parish before 1830 (newspaper advertisements and General Views of Agriculture); cereal suitability; share of heavy soils. All are taken unchanged from \citet{caprettinivoth2020data}. Parish area is the area of the 1851 parish polygon in the same data. Gonner's measure is the percentage of land enclosed by 1820 in the district; poor relief per head uses \texttt{POORRATE03}/\texttt{POP01}.

\textbf{Exposure.} For each parish, the years of all Acts linked to it. \textit{Recent}: any linked Act in 1790 to 1830. \textit{Old only}: linked Acts all before 1790. \textit{Future}: no linked Act before 1831 and a first linked Act in 1831 to 1870. Windows and years since the last Act are defined from the same years. Neighbour counts use parish centroids projected to kilometres at latitude 52.5\textdegree.

\textbf{Sample.} The main sample flag of the replication data ({{n.sample|int}} parishes) is used throughout; parishes with missing controls would be dropped, but none are missing in this sample. Standard errors are clustered on the {{n.counties|int}} counties (Yorkshire as three ridings); the 50 km grid has {{s.grid_se.clusters|int}} occupied cells.

{{TABLE_FULL}}

{{TABLE_LOR}}

{{TABLE_ROB}}

\end{document}
